\SourceLeaf{183}der völligen Unentschiedenheit in dieser Sache, falls jener
Beweis gegen die euklidische Geometrie entgegen meiner Erwartung nicht
als zwingend gelten sollte --- [sei von mir erkannt worden], teils weil
die vielen bisherigen Untersuchungen aus der ebenen Geometrie nicht
als verloren gelten müssten, sondern mit wenigen Änderungen brauchbar wären,
und weil wenigstens bis zu einer ziemlich weiten Grenze, die sich
vielleicht noch genauer bestimmen ließe, auch die Sätze der räumlichen
Geometrie und der Mechanik näherungsweise gültig wären, fand ich heute
Abend --- als ich gerade versuchte, einen Zugang zu Ihrer
transzendenten Trigonometrie zu finden, und weil es mir nicht gelang,
in der Ebene dafür hinreichende, bestimmte Funktionen zu erhalten,
zu räumlichen Konstruktionen überging, zu meiner nicht geringen
Überraschung folgenden \textls{Beweis für die euklidische Parallelentheorie}. .\ .\ .\ .\ .

.\ .\ .\ .\ .\ .\ Gerade als ich schließen will, bemerke ich noch, dass der obige
Beweis für die euklidische Parallelentheorie fehlerhaft ist. .\ .\ .\ .\ . Damit wäre
auch hier die Hoffnung verschwunden, zu einem völlig entschiedenen Ergebnis
zu kommen, und ich muss mich wieder mit dem oben Gesagten zufriedengeben.
Übrigens glaube ich, auf diesem Weg wenigstens einen Schritt zu Ihrer
transzendenten Trigonometrie getan zu haben, indem ich mithilfe der
sphärischen Trigonometrie die Verhältnisse aller Konstanten wenigstens
\textls{durch Konstruktion der rechtwinkligen} Dreiecke angeben kann. Es fehlt mir
noch die tatsächliche Berechnung der Basis eines gleichschenkligen Dreiecks aus
der Seite. Dafür werde ich versuchen, vom gleichseitigen Dreieck auszugehen.
.\ .\ .\}

\begin{center}\rule{3cm}{0.4pt}\end{center}

\Needspace{11\baselineskip}
\begin{center}\textsc{Bemerkung.}\end{center}

In dem Brief vom 12. Dezember 1816, aus dem hier Auszüge mitgeteilt werden,
bezieht sich \textsc{Wachter} wiederholt auf ein Gespräch mit \textsc{Gauss},
das während seines »letzten Aufenthalts in Göttingen« stattfand und von der
antieuklidischen Geometrie handelte. Aus einem Brief von
\textsc{Gauss} an \textsc{Olbers} vom 4. Juni 1816 geht hervor, dass dieser
Besuch \textsc{Wachter}s im April 1816 stattfand.

\begin{flushright}\textsc{Stäckel.}\end{flushright}

\begin{center}\rule{3cm}{0.4pt}\end{center}

